% ==============================================================================
% =================================== Aufgabe 1 =================================
% ==============================================================================

\chapter{Aufgabe}
\label{chap:1 Aufgabe}
\thispagestyle{fancy}
\section{} Wenn man moderne \uline{\acs{ERP}-Systeme} betrachtet, wird schnell deutlich, dass vor allem die \uline{hohe Wandlungsfähigkeit} über den langfristigen Erfolg entscheidet. Es geht dabei nicht nur um starre Software, sondern um die Fähigkeit, sich \uline{effizient} und \uline{schnell} an dynamische Marktanforderungen \uline{anzupassen}. In einem entsprechenden Bewertungsportfolio liegt der Idealbereich für zukunftsfähige Systeme genau dort, wie Gronau \cite[49]{1} in seiner Betrachtung anpassungsfähiger \acs{ERP}-Architekturen darlegt, wo sowohl eine \uline{hohe technische} als auch eine \uline{hohe geschäftsspezifische Wandlungsfähigkeit} vorliegt. \par Fünf entscheidende Erfolgsfaktoren für solche Systeme sind:
\begin{enumerate}[label=\arabic*.), leftmargin=*]
    \item \textbf{Hohe Wandlungsfähigkeit (\textit{Adaptivität}):} Ein System darf kein Klotz am Bein sein, sondern muss sich durch Skalierbarkeit und modularen Aufbau quasi mit dem Unternehmen mitentwickeln. Dies wird technisch durch Kriterien wie \uline{Skalierbarkeit}, \uline{Modularität} und \uline{Interoperabilität} sowie durch autopoietische Eigenschaften wie \uline{Selbstorganisation} unterstützt.
    \item \textbf{Integration von Künstlicher Intelligenz (\textit{\acs{KI}}) und Analytics:} Zukunftsfähige Systeme gehen über eine rein vergangenheitsorientierte Sicht hinaus und integrieren \uline{Monitoring-}, \uline{Prognose-}, \uline{Simulations-} und \uline{Optimierungsfunktionen}. Gronau \cite[340]{2} weist in seiner Architekturbetrachtung von \acs{ERP}‑Systemen darauf hin, dass die Nutzung von \acs{KI} eine unvoreingenommene Analyse vorhandener Daten ermöglicht, verlässliche Vorhersagen bietet und die Entscheidungsfindung durch digitale Assistenten unterstützt.
    \item \textbf{Rekonfigurierbare Geschäftsprozesse (\textit{"Modellieren statt Programmieren"}):} Aufgrund der zunehmenden \uline{Delinearisierung} von Arbeitsabläufen müssen zukünftige \acs{ERP}-Systeme Prozessmodelle aufweisen, die von den Nutzern zur Laufzeit individuell angepasst werden können. Die Einbindung von \uline{Workflow-Designern} ermöglicht es, Standardprozesse flexibel um spezialisierte Abläufe zu ergänzen.
    \item \textbf{Interoperabilität und Delokalisierung:} In einer vernetzten Welt müssen \acs{ERP}-Systeme die Unternehmensgrenzen öffnen und den Datenaustausch über die gesamte Wertschöpfungskette hinweg ermöglichen. Die Schnittstellen müssen \uline{interoperabel} und \uline{selbstbeschreibungsfähig} sein, um eine nahtlose Kommunikation zwischen unterschiedlichen Objekten des digitalen Unternehmens zu gewährleisten.
    \item \textbf{Skalierbare Infrastruktur und Cloud-Fähigkeit:} Mit Verweis auf die notwendige Leistungsfähigkeit für moderne Prozesse (\textit{ so wie transaktionale und analytische Prozesse}) unterstreicht Gronau (2021, S. 345–346) die essenzielle Rolle von \uline{Cloud-Architekturen}. Er argumentiert, dass diese Infrastrukturen notwendig sind, um die \uline{Skalierbarkeit des Geschäftsbetriebs} sowohl bei langfristigem (\textit{Wachstum}) als auch bei kurzfristigen saisonalen Schwankungen (\textit{Saisongeschäft}) sicherzustellen.
\end{enumerate}
\section{ } \par Für die innerbetriebliche Migration (\textit{Einführung}) eines \acs{ERP}-Systems setzt sich ein effektives \uline{Projektteam} aus verschiedenen hierarchischen und funktionalen Rollen zusammen:
\begin{itemize}
    \item Der \uline{Lenkungsausschuss} bildet das oberste Entscheidungsgremium und besteht idealerweise aus der Geschäftsführung, dem Gesamtprojektleiter, dem Projektleiter des Anbieters sowie einem \uline{Trusted Advisor}.
    \item Die Aufgaben dieses Ausschusses umfassen die Festlegung und Veränderung von Projektzielen, Terminen und Budgets sowie die Harmonisierung des Projekts mit den strategischen Unternehmenszielen.
    \item Im Rahmen der Maßnahmen zur Projektvorbereitung beschreibt Gronau \cite[308]{4} den \uline{Trusted Advisor} als einen Ratgeber der Geschäftsführung, der nicht nur die Qualität der Ergebnisse bewertet, sondern auch als kommunikative Brücke zwischen Anbieter und Kunde fungiert.
    \item Der \uline{Gesamtprojektleiter} trägt die operative Verantwortung, berichtet an den Lenkungsausschuss, genehmigt Projektphasen und überwacht die Kompatibilität zwischen der \acs{ERP}-Lösung und den Organisationsanforderungen.
    \item \uline{Key User} (\textit{Hauptanwender}) wirken als Multiplikatoren in ihren jeweiligen Fachbereichen, definieren Anforderungen in Workshops und leisten eine aktive Unterstützung bei der Konzeptumsetzung.
    \item Die eigentlichen \uline{Projektmitarbeiter} sollten zu mehr als 50\% ihrer Arbeitszeit für das Projekt freigestellt sein, da sie die Verantwortung für die Konzepterstellung und die Definition der Soll-Prozesse tragen.
    \item \uline{Externe Dienstleister} können als Berater oder Projektsteuerer ergänzend hinzugezogen werden, falls intern nicht ausreichend Know-how oder Kapazitäten verfügbar sind.
\end{itemize}
Die wichtigsten \uline{Ziele} des Migrationsprojekts lassen sich wie folgt beschreiben:
\begin{itemize}
    \item Bereits Gronau \cite[286]{5} weist darauf hin, dass eine klare \uline{Zieldefinition} zu Projektbeginn unerlässlich ist, um während der Laufzeit Orientierung zu bieten und nach Abschluss eine Erfolgskontrolle zu ermöglichen.
    \item Zu den primären \uline{organisatorischen Zielen} gehören die Verkürzung von Durchlaufzeiten, die Integration manueller Arbeitsabläufe und die Neugestaltung von Prozessen auf Basis von \uline{Best Practices}. Diese Ziele greifen die von Gronau \cite[287]{6} im Rahmen der Zieldefinition empfohlenen Maßnahmen auf.
    \item \uline{Wirtschaftliche Ziele} konzentrieren sich auf die Erreichung eines positiven \uline{Return on Investment (\textit{\acs{RoI}})} durch signifikante Kosteneinsparungen (\textit{\acs{z. B.} Verringerung des Materialaufwands}) und Produktivitätssteigerungen, so Gronau \cite[62]{7}.
    \item Des weiterem erklärt Gronau \cite[80]{8}, dass \uline{Technische Ziele} die Ablösung veralteter Software durch moderne, skalierbare Architekturen beinhalten sowie die Gewährleistung einer hohen Datenaktualität durch die Erfassung von Vorgängen in \uline{Echtzeit}.
    \item In Anlehnung an Gronau \cite[13]{9} und die dort zusammengestellten Vorteile von \acs{ERP}‑Systemen (\textit{\acs{vgl.} Sumner 2005, \acs{S.} 5}) besteht ein weiteres wesentliches Ziel darin, \uline{Transparenz} über sämtliche Unternehmensressourcen und Werteflüsse zu schaffen, um das Controlling und strategische Managemententscheidungen fundiert zu unterstützen.
    \item Zuletzt soll die langfristige \uline{Wandlungsfähigkeit} sichergestellt werden, damit das \acs{ERP}-System auch künftige Marktveränderungen und organisatorische Umbrüche effizient abbilden kann.
\end{itemize}
\section{ } \par Gronau \cite[253]{10} arbeitet fünf wesentliche Unterschiede und Gemeinsamkeiten in den Funktionen und Merkmalen von \acs{ERP}- und \acs{SCM}-Systemen heraus: \par \textbf{Unterschiede}
\begin{enumerate}[label=\arabic*.), leftmargin=*]
    \item \textbf{Fokus der Betrachtung:} Während \acs{ERP}-Systeme primär auf die \uline{unternehmensinterne Integration} und die Verwaltung interner Ressourcen ausgerichtet sind, fokussieren \acs{SCM}-Systeme auf die \uline{gesamte Lieferkette} und unternehmensübergreifende Wertschöpfungsnetzwerke.
    \item \textbf{Art der Planung:} Die Planung in \acs{ERP}-Systemen erfolgt meist \uline{sukzessiv} \acs{bzw.} \uline{sequenziell} nach dem \acs{MRP} II-Konzept, wohingegen \acs{SCM}-Systeme eine \uline{simultane Planung} (\textit{Advanced Planning and Scheduling – \acs{APS}}) ermöglichen.
    \item \textbf{Aktualität und Frequenz:} \acs{ERP}-Systeme arbeiten typischerweise mit täglichen oder wöchentlichen Intervallen, während \acs{SCM}-Systeme eine \uline{ständige Echtzeitverarbeitung} der Informationen erfordern.
    \item \textbf{Planungsgeschwindigkeit:} Ein Planungsvorgang in einem \acs{ERP}-System kann mehrere Stunden dauern, während \acs{SCM}-Systeme darauf ausgelegt sind, Planungsläufe in \uline{Minuten} abzuschließen.
    \item \textbf{Zielsetzung:} \acs{ERP}-Systeme dienen als Rückgrat der Informationsverarbeitung zur Administration interner Abläufe, während \acs{SCM}-Systeme spezifisch darauf zielen, netzwerkweite Probleme wie den \uline{Bullwhip-Effekt} durch Transparenz zu minimieren.
\end{enumerate}
\textbf{Gemeinsamkeiten}
\begin{enumerate}[label=\arabic*.), leftmargin=*]
    \item \textbf{Ressourcenmanagement:} Beide Systeme befassen sich grundlegend mit der \uline{Verwaltung} und \uline{Optimierung von Ressourcen}, wobei insbesondere der Materialfluss eine zentrale Rolle spielt.
    \item \textbf{Modularer Aufbau:} Wie Goram \cite[7]{11} unter Verweis auf Mertens (2007) erläutert, zeichnen sich sowohl \acs{ERP}- als auch \acs{SCM}-Systeme durch einen \uline{modulare Architektur} aus, der es Unternehmen ermöglicht, einzelne Funktionsbereiche je nach Bedarf einzusetzen.
    \item \textbf{Prozessintegration:} Beide Systemkategorien unterstützen die \uline{horizontale Integration} von Geschäftsprozessen. Dieses Merkmal greift Gronau \cite[243]{12} im Zusammenhang mit der Variantenbeherrschung auf und beschreibt es als zentralen Lösungsansatz, um den Informationsfluss entlang der gesamten Wertschöpfungskette zu verbessern.
    \item \textbf{Entscheidungsunterstützung:} Sie dienen als \uline{strategische Informationswerkzeuge}, die dem Management auf Basis von Datenanalysen und Berichten fundierte Entscheidungen ermöglichen. Diese Funktion hebt Gronau \cite[262]{13} insbesondere für die Distributionsplanung hervor.
    \item \textbf{Transparenz und Kostensenkung:} Ein wesentliches gemeinsames Ziel ist die \uline{Schaffung von Transparenz}, um Ineffizienzen (\textit{\acs{z. B.} Überbestände}) aufzudecken und Kosten nachhaltig zu senken. Goram \cite[9]{14} beschreibt dies als einen der zentralen Vorteile von \acs{ERP}-Systemen.
    \end{enumerate}

% ==============================================================================
% =================================== Aufgabe 1 =================================
% ==============================================================================